\subsection{Overall Benchmark Results}
\label{subsec:overall_benchmark_results}

This section evaluates the overall contribution of the proposed framework and then analyzes the VNS design choices that led to the selected configuration. The results are reported using relative gaps, best-instance counts, and average runtimes. For each benchmark instance \(i\), the relative gap of method \(m\) was computed with respect to the best objective value obtained among all compared methods:

\[
Gap_{im} =
\frac{Eval_{im} - Eval_i^{best}}
     {Eval_i^{best}}
\times 100 .
\]

Here, \(Eval_{im}\) denotes the objective value obtained by method \(m\) on instance \(i\), and \(Eval_i^{best}\) is the best objective value observed for that instance across the methods included in the relevant comparison. Average gap values reported in the tables correspond to the arithmetic mean of the instance-level gaps over the complete benchmark set. Therefore, lower average-gap values indicate better solution quality. The standard deviation measures the variability of these instance-level gaps, while the best-instance count reports the number of benchmark scenarios in which a method achieved the best observed objective value.

Table~\ref{tab:main_benchmark_comparison} compares the selected proposed configuration against the principal benchmark variants. The selected configuration corresponds to the adaptive \(0.20\)--\(0.30\) VNS policy with \(K_{\max}=200\). It is used as the reference method in this comparison because it achieved the best observed objective value for every benchmark scenario.

\begin{table}[H]
\centering
\caption{Overall comparison of the selected proposed configuration and the main benchmark variants.}
\label{tab:main_benchmark_comparison}
\begin{tabular}{lcccc}
\hline
\textbf{Method}
& \textbf{Average gap (\%)}
& \textbf{Std. dev. (\%)}
& \textbf{Best-instance count}
& \textbf{Average runtime (min)} \\
\hline
Fixed sample & 361.93 & 198.43 & 0 & 0.1 \\
No outer reconfiguration & 140.07 & 90.32 & 0 & 17.1 \\
No VNS diversification & 178.67 & 113.52 & 0 & 0.5 \\
Selected proposed configuration & 0.00 & 0.00 & 90 & 85.1 \\
\hline
\end{tabular}
\end{table}

The selected proposed configuration dominates the main benchmark variants in terms of solution quality. Most importantly, it achieved the best solution in every one of the \(90\) benchmark scenarios, whereas none of the competing variants achieved a single best-instance result. This is the strongest aggregate result in the benchmark study because it shows that the improvement is not driven by a small subset of favorable instances. Instead, the selected configuration is consistently superior across all combinations of due-date tightness, test portfolios, voltage requirements, and workload conditions represented in the benchmark set.

The fixed-sample policy performs worst, with an average gap of \(361.93\%\) and a standard deviation of \(198.43\%\). Although this policy is computationally very fast, its poor performance indicates that the industrial rule of assigning a fixed number of samples to every product is not robust under heterogeneous planning conditions. The fixed rule ignores the fact that the marginal value of an additional sample differs across products, test environments, and due-date settings. As a result, it may allocate testing effort to products for which additional samples provide little scheduling benefit while failing to relieve bottlenecks for more critical projects.

The no-outer-reconfiguration variant also performs substantially worse than the selected proposed configuration, with an average gap of \(140.07\%\). This result demonstrates that sample-size optimization alone is insufficient when chamber settings are not repeatedly updated. Changes in the sample vector alter environment-specific workload distributions, voltage requirements, and chamber utilization patterns. If the chamber configuration remains fixed while the workload structure changes, the algorithm evaluates sample decisions under a resource structure that may no longer be appropriate.

The no-VNS-diversification variant has an average gap of \(178.67\%\), confirming the importance of diversification after local search and chamber reconfiguration have converged. Without VNS, the search can remain trapped in workload structures that cannot be improved through deterministic local moves alone. Its shorter runtime therefore reflects the absence of a critical improvement mechanism rather than an efficient path to high-quality solutions.

\subsection{VNS Ratio Analysis}
\label{subsec:vns_ratio_analysis}

After establishing the overall advantage of the selected configuration, the VNS design was analyzed in more detail. Table~\ref{tab:vns_ratio_comparison} compares alternative VNS perturbation-ratio policies under a fixed shake budget of \(K_{\max}=100\). This comparison isolates the effect of perturbation intensity on the balance between diversification and preservation of useful solution structure.

\begin{table}[H]
\centering
\caption{Comparison of VNS perturbation-ratio policies with \(K_{\max}=100\).}
\label{tab:vns_ratio_comparison}
\small
\begin{tabular}{lccccc}
\hline
\textbf{Measure}
&
\shortstack{\textbf{0.10}\\\textbf{VNS (100)}}
&
\shortstack{\textbf{0.20}\\\textbf{VNS (100)}}
&
\shortstack{\textbf{0.30}\\\textbf{VNS (100)}}
&
\shortstack{\textbf{0.10--0.20--0.30}\\\textbf{VNS (100)}}
&
\shortstack{\textbf{0.20--0.30}\\\textbf{VNS (100)}} \\
\hline
Average gap (\%) & 29.62 & 18.11 & 23.44 & 22.06 & 14.93 \\
Std. dev. (\%) & 44.09 & 31.52 & 30.10 & 33.25 & 27.35 \\
Best-instance count & 25 & 27 & 18 & 38 & 36 \\
Average runtime (min) & 41.3 & 32.1 & 30.4 & 66.0 & 49.7 \\
\hline
\end{tabular}
\end{table}

The results show that a single perturbation ratio is not sufficient to obtain consistently strong performance. Among the single-ratio variants, \(0.20\) performs best on average, with an average gap of \(18.11\%\). By contrast, \(0.10\) produces a larger average gap of \(29.62\%\), suggesting that very mild perturbations often fail to move the search out of the basin of attraction reached by the local improvement phase. The \(0.30\) ratio performs better than \(0.10\), but its average gap remains \(23.44\%\), which indicates that stronger perturbations alone can be too disruptive. They create diversity, but may also destroy useful sample-allocation and chamber-workload structures already constructed by the preceding search phases.

The adaptive fallback policies improve robustness by applying perturbations sequentially instead of committing to a single disruption level. The \(0.20\)--\(0.30\) policy achieves the lowest average gap among the \(K_{\max}=100\) variants, with an average gap of \(14.93\%\), and also has the lowest standard deviation in this group. This behavior is consistent with the intended role of the fallback mechanism: the search first attempts to improve the solution through a moderate perturbation, and only escalates to a stronger move when the moderate shake is insufficient.

Although the \(0.10\)--\(0.20\)--\(0.30\) policy achieved the highest number of best-instance results, its average gap and runtime were considerably worse. Since the objective of tactical planning is to obtain robust performance over the entire benchmark set rather than maximizing the number of isolated wins, average performance and variability were given higher priority than best-instance count. For this reason, the \(0.20\)--\(0.30\) fallback structure was selected as the preferred VNS ratio policy.

\subsection{Shake Budget Analysis}
\label{subsec:shake_budget_analysis}

After selecting the \(0.20\)--\(0.30\) fallback policy, the effect of the shake budget was evaluated by comparing \(K_{\max}=50\), \(100\), and \(200\). The shake budget controls how much opportunity the VNS phase has to generate and evaluate alternative sample structures after local improvement stagnates. The results are summarized in Table~\ref{tab:vns_budget_comparison}.

\begin{table}[H]
\centering
\caption{Effect of shake budget for the \(0.20\)--\(0.30\) VNS policy.}
\label{tab:vns_budget_comparison}
\begin{tabular}{lccc}
\hline
\textbf{Measure}
& \textbf{0.20--0.30 VNS (50)}
& \textbf{0.20--0.30 VNS (100)}
& \textbf{0.20--0.30 VNS (200)} \\
\hline
Average gap (\%) & 13.79 & 5.65 & 0.31 \\
Std. dev. (\%) & 29.07 & 20.06 & 1.65 \\
Best-instance count & 18 & 47 & 82 \\
Average runtime (min) & 32.6 & 49.7 & 85.1 \\
\hline
\end{tabular}
\end{table}

Increasing the shake budget clearly improves both average performance and robustness. With \(K_{\max}=50\), the average gap is \(13.79\%\), indicating that the diversification phase is often terminated before sufficient alternative workload structures are explored. Increasing the budget to \(100\) reduces the average gap to \(5.65\%\) and increases the number of best-instance results from \(18\) to \(47\). The improvement is not limited to a few favorable cases: the standard deviation decreases from \(29.07\%\) to \(20.06\%\), showing that the additional shake attempts also reduce performance variability.

The best overall performance is obtained with \(K_{\max}=200\). This setting reduces the average gap to \(0.31\%\), decreases the standard deviation to \(1.65\%\), and achieves the best result in \(82\) out of \(90\) scenario combinations within the shake-budget comparison. The increase from \(K_{\max}=100\) to \(K_{\max}=200\) reduced the average gap from \(5.65\%\) to \(0.31\%\) and increased the best-instance count from \(47\) to \(82\). Such improvements are substantial and justify the corresponding increase in runtime for the tactical planning context considered in this study.

The average runtime increases from \(32.6\) minutes for \(K_{\max}=50\) to \(49.7\) minutes for \(K_{\max}=100\), and to \(85.1\) minutes for \(K_{\max}=200\). This additional computational effort is acceptable because the framework is intended for tactical R\&D test planning, where solution quality and robustness are more important than real-time response. The \(K_{\max}=200\) setting was therefore selected because it provides a much more reliable final configuration, not merely because it produces the lowest average gap.

\subsection{Managerial Insights}
\label{subsec:managerial_insights}

The benchmark results provide several managerial insights for tactical test planning. First, sample counts should be treated as decision variables rather than fixed operational inputs. The poor performance of the fixed-sample policy shows that a uniform sampling rule cannot respond effectively to heterogeneous due dates, test portfolios, voltage requirements, and chamber bottlenecks. Flexible sample-size decisions allow the planning system to allocate testing effort where it has the greatest impact on tardiness reduction.

Second, sample-size optimization and chamber reconfiguration should be treated as coupled decisions. The no-outer-reconfiguration variant demonstrates that improving sample counts without updating chamber settings leaves a large performance gap. In practical terms, this means that a change in project-level sample decisions should trigger a reassessment of the chamber configuration, because the revised workload may shift demand across environments and voltage categories.

Third, diversification is essential even after local improvement has converged. The no-VNS-diversification variant is faster, but it does not produce a best solution in any benchmark scenario. This indicates that deterministic local moves alone are not sufficient for the complex workload structures generated by the test-planning problem. Controlled VNS perturbations provide a mechanism for escaping stagnant regions while preserving useful structure from high-quality incumbent solutions.

Overall, the selected proposed configuration reflects a deliberate quality--runtime trade-off. The adaptive \(0.20\)--\(0.30\) VNS policy avoids both insufficient and excessive perturbation, while \(K_{\max}=200\) provides enough diversification effort to make the results stable across the benchmark set. The evidence supports the main premise of the study: effective tactical test planning requires joint optimization of sample counts, chamber configuration, and diversification strategy rather than treating these decisions independently.
